\section{Canonical whole tensor-power pullback coherence}

\begin{lemma}\label{lem:actual-power-iso-functoriality}
\lean{EnvelopeProperty.BundleTensorPowerPullbackCoherence.powerIso_refl,EnvelopeProperty.BundleTensorPowerPullbackCoherence.powerIso_trans,EnvelopeProperty.BundleTensorPowerPullbackCoherence.powerIso_symm}
\leanok
\uses{lem:bundle-tensor-power-isomorphism,lem:cartier-actual-monoidal-explicit-tensor-naturality}
The genuine explicit tensor-power construction preserves identity,
composition and inverse of every whole original module isomorphism.
\end{lemma}
\begin{proof}
The explicit tensor of two isomorphisms is the tensor of their
underlying module-presheaf isomorphisms followed by sheafification.
Both operations preserve identity and composition: for composition,
the tensor interchange law combines the two maps in each factor,
and functoriality of sheafification combines their images. Induct
on the power. At zero the power isomorphism is the identity of the
structure sheaf. At a successor the preceding tensor identities
and the induction hypothesis prove identity and composition.
Apply the composition identity to an isomorphism and its inverse;
the identity case then identifies the power of the inverse with
the inverse power isomorphism.
\end{proof}

\begin{lemma}\label{lem:actual-power-pullback-successor}
\lean{EnvelopeProperty.BundleTensorPowerPullbackCoherence.pullbackIso_succ}
\leanok
\uses{lem:bundle-tensor-power-pullback,lem:cartier-actual-monoidal-explicit-tensor-naturality}
For every actual $f:Y\to X$ and module $M$, the canonical comparison
$f^*(M^{\otimes(n+1)})\simeq(f^*M)^{\otimes(n+1)}$ is the genuine
explicit tensor pullback comparison followed by the tensor of the
previous power comparison and the identity of $f^*M$.
\end{lemma}
\begin{proof}
Expand the actual power pullback comparison. Naturality of the
native-to-explicit tensor comparison moves the tensor of the
previous power comparison and the identity past the final
explicit comparison. The remaining three maps are exactly the
canonical explicit tensor pullback comparison. Associativity of
composition gives the asserted equality of whole isomorphisms.
\end{proof}

\begin{lemma}\label{lem:actual-power-pullback-naturality}
\lean{EnvelopeProperty.BundleTensorPowerPullbackCoherence.pullbackIso_naturality}
\leanok
\uses{lem:actual-power-pullback-successor,lem:cartier-original-tensor-pullback-naturality}
The genuine tensor-power pullback comparisons are natural in every
whole original module isomorphism, for every nonnegative power.
\end{lemma}
\begin{proof}
At zero both power isomorphisms are identity maps, so the two
composites are the native structure-unit comparison. At a successor
expand the preceding recursion on both sides. Naturality of the
canonical explicit tensor pullback comparison moves the tensor of
the original power isomorphism and the original isomorphism through
that comparison. Tensor functoriality then combines this map with
the two pullback power comparisons. The induction hypothesis in
the first factor and the unchanged map in the second factor give
exactly the other composite.
\end{proof}

\begin{theorem}\label{thm:actual-power-pullback-composition}
\lean{EnvelopeProperty.BundleTensorPowerPullbackCoherence.pullbackIso_comp,EnvelopeProperty.BundleTensorPowerPullbackCoherence.pullbackIso_congr}
\leanok
\uses{lem:actual-power-iso-functoriality,lem:actual-power-pullback-successor,lem:cartier-canonical-explicit-tensor-pullback-composition,lem:actual-power-pullback-naturality}
For actual composable scheme maps, the genuine whole tensor-power
pullback comparison commutes with the native pullback-composition
isomorphism, including its power on the original module. It also
commutes with the native comparison of any two equal scheme maps.
\end{theorem}
\begin{proof}
For composition induct on the power. At zero the statement is the
native composition identity for the structure-unit comparisons.
At a successor expand both power comparisons using the explicit
tensor recursion. First apply naturality of explicit tensor
pullback to the previous power comparison. The tensor interchange
law then combines the previous two power comparisons with the
power of the native composition isomorphism. The induction
hypothesis identifies the first factor; the second factor is the
native composition isomorphism itself. The genuine composition
identity for canonical explicit tensor pullback identifies the
remaining composite with the native composition comparison on
the original successor tensor power followed by its direct
pullback comparison. This is the claimed whole isomorphism.
For equal maps substitute the equality. The native equal-map
comparison becomes identity, and its power becomes identity by
the preceding functoriality lemma, proving that identity as well.
\end{proof}
